% R. Brown's translations of two letters to L. Breen, with his note and comments (cote 134-2): archive cote 134-2, site page 118 (PDF page 119).
% comments. Draft compared: none.
% AI-transcribed by Claude Opus 5.5 (claude-opus-5-5, adaptive thinking, effort high), Message Batches API, batch pilotnext (msgbatch_01A9sdu8SYq23bEtqX9tUncQ), 2026-11-13; not yet checked by eye.
\scanpage{134-2/118}
\mspage{-2-}

equivalences between six categories as set out in the diagram,

\diagram{\begin{tikzcd}[column sep=large, row sep=large]
\text{poly -T-complexes} \arrow[r, leftrightarrow, "(f)"] \arrow[d, leftrightarrow, "(f)"'] & \text{cubical T-complexes} \arrow[r, leftrightarrow, "(c)"] & \omega\text{-groupoids} \arrow[d, leftrightarrow, "(d)"] \arrow[dl, leftrightarrow, "(a)"] \\
\text{simplicial T-complexes} \arrow[r, leftrightarrow, "(e)"'] & \text{crossed complexes} \arrow[r, leftrightarrow, "(d)"'] & \infty\text{-groupoids}
\end{tikzcd}}

The categories are defined and the equivalences proved in:

a) R. Brown and P. J. Higgins, 'The algebra of cubes',
J. Pure Appl. Alg. 21(1981) 233 - 260.

b) R. Brown and P. J. Higgins, 'Colimit theorems for relative
homotopy groups', \textit{ibid.} 22(1981) 11 - 41.

c) R. Brown and P. J. Higgins, 'The equivalence of $\omega$-groupoids and
cubical T-complexes', Cah. Top. Géom. Diff. (3$^{\text{e}}$ Coll. sur les
catégories, dédié a Charles Ehresmann) 22 (1981), 349 - 369.

d) R. Brown and P. J. Higgins, 'The equivalence of $\infty$-groupoids and
crossed complexes' \textit{ibid,} 370 - 386.

e) N. K. Ashley, \textit{Simplicial T-complexes,} Ph.D. thesis, University
of Wales (1978),(to be published in an issue of Esquisses
Mathématiques, Amiens).

f) D. W. Jones, Ph.D. thesis in preparation.

Applications of the equivalence (a) are given in reference (b)
above.

5. However the notion of $\infty$-groupoid referred to in (4) is not adequate
for the purposes considered by Grothendieck since, roughly speaking,
crossed complexes describe the homotopy types only of spaces X
which fibre over a K($\pi$,1) with fibre a product of Eilenberg-
MacLane spaces.

6. A notion of multiple groupoid called an \textit{n-cat-group} has been described

\mspage{\hand{119}}
